\section{公式揭示的新物理}

复现的落脚点不是「公式算对了」，而是这些精确公式揭示了长波近似掩盖的物理。四个洞察：

\subsection{长波近似的失效边界}

Table 1（长波近似）把球 Bessel 核截断为 $j_l(kr)\approx(kr)^l/(2l+1)!!$，只在小尺寸参数成立。精确的 Table 2 保留完整核，两者的差异在磁共振附近变得不可忽略：

\begin{itemize}
  \item 介电球 $s=0.75$ 处，长波近似给出 ED $136\%$、MD $277\%$ 的误差（$>100\%$）；
  \item 精确 Table 2 与 Mie 一致到 $<0.21\%$。
\end{itemize}

这揭示：多极矩不是静态的 $\int\mathbf r\times\mathbf J$，而是带波长依赖振荡核的积分——纳米结构自身的尺寸（相对波长）参与了多极的定义。长波近似只在 $kr\ll1$ 可用，磁共振区（$kr\sim1$）失效。

\subsection{toroidal 多极不是独立自由度}

FC2017 严格证明：把电偶极的精确展开按电部分与 toroidal 部分拆分后，两部分各自含不可观测的非辐射分量，求和才抵消。因此 toroidal 无法被单独测量，它只是电偶极系数对电磁尺寸展开的次低阶修正项（Table 1 的 $k^2/10$ 项）。这修正了「多极家族含独立 toroidal 成员」的图像——toroidal 是描述同一电偶极的另一种展开方式，不是新自由度。

\subsection{金属纳米结构的磁响应}

双金盘（$a=250$ nm）在 $x=0.36$（$\lambda=1389$ nm）处磁偶极主导（MD $1.60\times10^{-12}$ m$^2$），电四极伴生、磁四极近零。这说明光学频段的金属结构（不只是介质）有显著的磁偶极共振响应，而长波近似会低估这一磁贡献。精确多极分解才能正确分辨电/磁各阶的贡献。

\subsection{多极干涉与 Fano 起源}

双金盘的方向交叉项从 $-24.9$（$x=0.82$，相消）到 $+4.5$（$x=0.53$，相长），明确区分了多极叠加的相消与相长。这正是 Fano 共振的物理起源——不是单个多极的行为，而是电偶极与更高阶多极之间的干涉。
